\section{Conjugate alternation test}\label{sec:alternation-imaginary}

The parity needed for residual concentration is a consequence of an alternating Bockstein pairing. We first prove that assertion integrally, including at quadratic ramification primes, and then apply it to the imaginary \(S_3\) case of Proposition~\ref{prop:anchor}. Throughout, the finite-precision diagrams and compatible limits are those of Section~\ref{sec:foundations}; all comparisons are made at the stages before passing to a limit. Over an imaginary quadratic \(K\) consider \(T\chi\), \(T\) a free plane from \(\mathbb Q\) with a perfect alternating determinant form \(e:T\otimes T\to R(1)\), \(\chi\) an anticyclotomic character. Here \(R\) can be our complete scalar integer or order/power-series ring; Galois actions, including their quadratic-base compatibilities and the form, come from actual stage actions at finite precisions. Allow the quadratic ramification primes in inducing. Write
\[
 V=\operatorname{Ind}_K^{\mathbb Q}(T\chi)=V_1\oplus V_2,\qquad
 B((a,b),(a',b'))=e(a,b')+e(b,a').
\]
The copies over \(K\) use opposite scalar twists; a fixed complex conjugation \(c\) interchanges them with the ordinary \(T\)-action. This is the conjugate Tate pairing; Shapiro pushes the corresponding cup by trace to the compact scalar invariant over \(\mathbb Q\), compatible as in the diagram discussion (no factor two). The real Tate complex for \(V\) is contractible integrally.

Use the following local conditions for \(V\):
\begin{enumerate}[label=\textup{(\roman*)}]
\item One unrestricted coordinate at a split prime pair, with the other coordinate strict; a zero local condition is also allowed.
\item At each newly allowed odd discriminant prime \(q\), assume that \(T\) is good and unramified and that \(\chi\) is unramified over \(K\), and impose the unramified cochains over \(K\).
\end{enumerate}

For the latter one can equivalently use residue cochains over \(\mathbb Q_q\) on the \(V\)-inertia invariants (evaluation by Shapiro on the original summand, same residue field). In fact quadratic inertia acts by a swap \((a,b)\mapsto(\delta b,\delta^{-1}a)\); the invariants are the graph \(b=\delta^{-1}a\), with \(\delta\) a character unit, compatible by reduction. Cup here has the unramified isotropy trivialization (residue cochains of 2-cohomological dimension one); on a split coordinate cup is identically zero on the condition. These give the usual compact cup on the Selmer complex. Split coordinates and the stated unramified conditions are exact self-orthogonals: check locally also by Shapiro over \(K\), for the latter using residue-field-coefficient duality, unramified inclusion of all degree zero and injection in degree one, and the odd local Euler formula (unramified \(H^1\) has dimension \(h^0\), complementary to the dimension of the dual unramified space). Then use finite models to lift the orthogonality isomorphism. Zero conditions need not be self-orthogonal in using the test by itself.

\begin{lemma}[Conjugate Bockstein alternation]\label{alternation-imaginary:alternation}
With the coefficient module, local conditions, and isotropy data just specified, let \(s\in R\) be a regular scalar. On degree one for coefficients mod \(s\), with \(\partial\) the Bockstein from mod \(s^2\), one has
\[
                           \langle x,\partial x\rangle=0 .             \tag{H1}\label{eq:H1}
\]
\end{lemma}
\begin{proof}
\emph{The central extension.}
On \(V_{s^2}\times R_{s^2}(1)\), with subscripts denoting reductions, use the law
\[
 (a,b,z)(a',b',z')
  =(a+a',b+b',z+z'+e(a,b')).
\]
Its commutator is \(B\). The actions over \(K\) are diagonal as above. If \(C\) denotes the action of complex conjugation on \(T\), then \(e(Ca,Cb)=-e(a,b)\), and the lift of \(c\) is
\[
 (a,b,z)\longmapsto(Cb,Ca,-z+e(a,b)).
\]
Equivalently, its central coordinate is \(c(z-e(a,b))\). Alternation of \(e\) shows directly that this preserves the group law, is an involution, and respects conjugation of the \(G_K\)-actions. On the vectors \(sV_s\) the lift \(y\mapsto(sy,0)\) is now additively and equivariantly split; conjugation by a lift with underlying vector \(v\) adds \(sB(v,y)\) to the central coordinate.

\emph{Global and local defects.}
Write \(x\) also for a global crossed cocycle. Take a cochain lift to this group, with a linear-coordinate lift as in computing the Bockstein. Its product defect \(\widetilde x(g)g\widetilde x(h)\widetilde x(gh)^{-1}=(s y(g,h),\delta_x(g,h))\) uses the global Bockstein \(y\); associativity gives \(d\delta_x=-sB(x\cup y)\). Write locally \(x_v=X+d a_v\) with \(X\) the cocycle in the prescribed local condition (at the real place use zero and ordinary positive-degree representatives). Make a reference lift \(A_v^{-1}\widetilde X(g)g A_v\), using lifts of \(a_v,X\) from the linear Bockstein calculation, the latter inside the allowed vectors with central coordinate zero. In the unramified case use the lifted invariant graph; inertia is trivial on the graph-with-center since \(e(a,\delta^{-1}a)=0\) and the place is odd. If the defect of \(\widetilde X\) is \((sY,\delta_U)\), the reference-lift defect is \((sY,\delta_U-sB(a_v,Y))\). Here \(\delta_U\) is zero for coordinate or zero conditions, unramified in the graph case. Write the actual localized global lift as \((s b_v,\ell_v)\) times the reference lift. Direct multiplication gives
\[
                y_v=Y+d b_v,\qquad
   \delta_{x,v}=\delta_U-sB(a_v,Y)+d\ell_v+sB(x_v\cup b_v),
\]
where \(b_v\) on the linear side denotes the Bockstein localization homotopy. Choose a primitive \(p_0\) for \(B(X\cup Y)\), zero on split/zero conditions and in residue cochains for the unramified condition (there its degree-two ambiguity is exact). Then \(\delta_U+s p_0\) is locally exact, again by residue cohomological dimension in the latter case. The compact pairing uses \(B(x\cup y)\) with local primitives
\[
                     p_0+B(a_v\cup Y)-B(x_v\cup b_v)
\]
by the mapping fibers with these isotropy data. At infinity the ordinary local computation supplies the needed Tate primitive relations in positive degrees. The formulas show that multiplying this compact class by \(s\) into the mod-\(s^2\) compact cochains gives zero. Compact invariant is scalar-compatible, and \(R/(s)\to R/(s^2)\) by \(s\) is injective, proving \eqref{eq:H1}.

\emph{Finite-precision passage to limits.}
For limit systems in this calculation, transfer the Shapiro models, localization triples, compact cup and trace as in the finite-model lemma, including unramified null via the residue cochains. Choose in limit free models the linear lifts of the Selmer complex data modulo \(s^2\) with differentials \(s\) times the chosen Bockstein data. Use every fixed Artin precision also of \(R/(s^2)\) and its further mod-\(s\) reduction. The equations in finite models then hold over those rings on stage sets in the ultrafilter. Transferring back by the compatible cochain maps gives actual local and global linear cochains at that precision, in particular \(d\widetilde x_{\rm lin}=s y,\ \widetilde x_{{\rm lin},v}-\widetilde X_{\rm lin}-d\widetilde a_v=s b_v\), with the indicated cochain equations already satisfied modulo \(s\). Graph/residue or coordinate cochains use their true lifts and maps in this procedure. The calculation killing \(s\) times the compact cup therefore applies even if this precision has a further annihilator for \(s\): it multiplies using the square-zero image ideal and these linear equations. Passage to the full reductions proves \eqref{eq:H1}. At places where both a more general local condition and the full local complex become contractible after a localization one can replace that condition by zero for the localized test. The pairings restrict compatibly (the zero restriction needs no isotropy correction); lift from the replaced integral complex mod \(s\) after clearing a scalar unit of the localization and invert again. Likewise one may base change first to an integral power-series curve as in our specialization conventions.

\end{proof}

\begin{corollary}[Parity of special and generic ranks]\label{alternation-imaginary:parity}
For a perfect two-term DVR complex in degrees \(1,2\) with perfect self-duality of shift \(-3\), arising from the preceding local conditions and satisfying \eqref{eq:H1}, the generic and special first-cohomology dimensions have the same parity.
\end{corollary}
\begin{proof}
Choose a uniformizer \(\varpi\) and \(s=\varpi^N\), where \(N\) is at least the length of every finite elementary divisor. The allowed scalar base changes permit this choice. Modulo \(s\), duality between \(H^1\) and \(H^2\) is perfect, since the principal Artin ring is self-injective. For a differential block \(\varpi^a\), with \(a\le N\), the degree-one group is \(\varpi^{N-a}R/\varpi^N\), the degree-two group is \(R/\varpi^a\), and the Bockstein sends the class of \(\varpi^{N-a}\) to the class of \(1\). It is therefore an isomorphism on the two groups of this block, whereas it vanishes on free-lifting classes. Thus the alternating Bockstein form on \(H^1\) has radical given by the free-lifting part and gives a perfect alternating pairing on the quotient (injectivity there and length equality with its dual). Splitting hyperbolic pairs of maximal order shows the number of finite nonunit divisors is even, as required.
\end{proof}

\section{The imaginary \texorpdfstring{\(S_3\)}{S3} anchor}\label{alternation-imaginary:anchor-section}

\begin{proposition}[Imaginary \(S_3\) anchor]\label{alternation-imaginary:anchor}
Proposition~\ref{prop:anchor} holds for every non-CM elliptic curve \(E/\mathbb Q\) whose residual image is \(S_3\) and whose quadratic residual subfield is imaginary.
\end{proposition}

The proof has three separate tasks: choose a CM comparator with known simple analytic product, arrange ordinary Selmer corank one for a twist of \(E\), and transfer an integral determinant unit from the comparator. The Selmer seed alone is not used to infer analytic rank.

\Needspace{12\baselineskip}
\subsection{The comparator and the ordinary Selmer seed}

\begin{lemma}[Choice of comparator]\label{alternation-imaginary:comparator}
Let \(W=E[2]\) have image \(S_3\) and imaginary quadratic subfield \(L\). There exist a primitive CM newform \(f_0\) of weight two and trivial nebentypus, lifting \(W\) on the plane \(M^*\) of Section~\ref{sec:cm}, and an imaginary quadratic field \(K\) of odd fundamental discriminant \(k<-4\), such that
\[
 L(f_0,1)\ne0,\qquad \ord_{s=1}L(f_0/K,s)=1,
 \qquad q\mid 2N_E N_{f_0}\Longrightarrow q\text{ splits in }K.
\]
Moreover, \(K\ne L\), and \(W|_{G_K}\) still has image \(S_3\).
\end{lemma}
\begin{proof}

Suppose \(E\) is non-CM with \(W=E[2]\) having image \(S_3\) and imaginary quadratic subfield \(L\). Use a CM newform \(f_0\) lifting \(W\) on \(M^*\) as constructed in the CM comparison. Quadratic twists retain this lifting property; take \(L(f_0,1)\ne0\), then take imaginary \(K\) of odd fundamental discriminant \(k<-4\) splitting \(2N_E N_{f_0}\) with \(L(f_0/K,s)\) having a simple zero at 1. These choices use the nonvanishing theorems in Section~\ref{sec:classical}, with their compatible sign and local splitting prescriptions \cite{FH95,BFH90}. If desired, the comparator may first be twisted by a quadratic character ramified at every prime of \(2N_E\) where the form's original level was good, locally trivial at primes already dividing that level (squareclass approximation). Those good local representations become ramified. Use a subsequent coprime twist split at that level with archimedean sign chosen for functional sign plus, for value nonvanishing. Thus all splitting primes required for derivative nonvanishing can be included in the actual level. Use the primitive form after twisting. Since \(L\) ramifies within its level, \(K\ne L\), \(W|_{G_K}\) still has full image.

\end{proof}

\begin{lemma}[Simultaneous Selmer seed]\label{alternation-imaginary:selmer-seed}
For the field \(K\) above, there is a squarefree product \(h\) of fresh primes congruent to \(1\pmod8\), coprime to \(2N_Ek\) and split in \(K\), such that
\[
 \dim_{\mathbb F_2}\Sel_2(E^h/\mathbb Q)
 +\dim_{\mathbb F_2}\Sel_2(E^{hk}/\mathbb Q)=1,
 \qquad s_2(E^h/K)=1.
\]
\end{lemma}
\begin{proof}
We construct \(h\) by successive fresh twists. The two curves \(E,E^k\) have opposite signs and no rational two-torsion, hence opposite finite two-Selmer dimension parities by 2-parity and Cassels alternation. Twist each time by a fresh positive good prime local square at all old relevant places, splitting in \(K\) and on \(W\); signs and conditions away from the new prime are unchanged. There on both curves the unramified plane switches to a transverse plane (Kummers of invariant two-torsion on the ramified twist evaluate surjectively on inertia, untwisted halves unramified and twist inertia negating them). By reciprocity and exact Kummer orthogonality the dimension change is \(\le 2-2j\) on either curve for old evaluation rank \(j\) there: new singular contributions annihilate the finite evaluations and the kernel is old strict there. Frobenius evaluations can be prescribed on a basis of the span of both Selmers in \(H^1(\mathbb Q,W)\): restriction to the joint kernel of \(W\) and abelian congruence data detects classes by the normal \(C_3\) without invariants, and additive joint evaluations span by absolute simplicity (the ground field is just \(\mathbb F_2\)). Use Chebotarev; the simultaneous residual evaluation statement is also the one in \cite[Lemma~3.5]{MazurRubin10}. First evaluate with rank two on the odd group until it has dimension one. Then whenever the even group has dimension at least two evaluate with rank two on it and rank one on the odd group simultaneously (possible even for a line inside the even group). The resulting dimensions \(0,1\) and quadratic isogeny with parity give the assertion.
\end{proof}

In the remainder of this section, \(E\) denotes the replacement \(E^h\). The product \(h\) will be restored when we conclude the proof of Proposition~\ref{alternation-imaginary:anchor}.

\subsection{Paired measures and strict complexes}

Use both \(T=T_2 E\otimes\mathcal O\) and \(T=M^*\) for \(f_0\) over common dyadic integers \(\mathcal O\) with uniformizer \(\pi\), and common split rational support \(S_f\) containing 2 and both level supports. For \(M^*\) use the unimodular alternating form \(e\) normalized by \(\mu^{-1}\) as at \eqref{eq:C1}, and use \(\lambda_A\) in sending Heegner Kummers on \(A=e_{f_0}J\) to dual Tate projected to this plane (\(\log_\omega\) on them incorporates \(\lambda_A^{-1}\), with \(\omega\) as there). For each plane use the strict complex (strict only at \(w\mid2\) in \(K\), full at all other finite allowed places), and paired CM-disk measures with character variables:
\begin{enumerate}[label=\textup{(\roman*)}]
\item \(1+t\) is the free anticyclotomic direction of the orders of 2-power conductor in \(K\).
\item The finitely many variables \(1+u_j\) come from cyclic quotients for varying inert primes \(r_j\), which are also allowed in the support. Their order conductors are \(r_j\), their characters are primitive on relative inertia, and their quotient orders \(2^{m_j}\) tend to infinity.
\end{enumerate}

All measure and arithmetic limits are on common stages; write \(R=\mathcal O[[t,\mathbf u]]\) on the arithmetic side and use the unramified constant limit with fixed further scalars as before for measures. The strict complex with total scalar twist has square amplitude 1,2 (residual invariant vanishing over \(K\), Poitou--Tate and global and dyadic Euler characteristics); denote its determinant by \(D\). Denote the measure product by \(B=b^+b^-\), depleting at exactly \(S_f\). The letters \(D,B\) refer separately to the two coefficient planes; when comparing them we write \(D_E,B_E\) and \(D_0,B_0\). Indeed \eqref{eq:M1}--\eqref{eq:M3} apply as described there with odd conductor \(s=\prod_j r_j\); the 2-part from the split order class sequence uses the unit ratios with image in the free direction of fixed finite index, compatibly through the odd conductors. Use common underlying disk/class and Serre--Tate data and compatible tame orientations across levels, canonical connected level at 2. At conductor \(s\) take basepoints by simultaneous descending cyclic quotients from a conductor-one split point, transporting the level. The fully depleted Tate expansions, their integral congruence test, and the common residual Hecke data \(W\) give
\[
 B_E\equiv B_0\pmod\pi,\qquad
 D_E\equiv\epsilon D_0\pmod\pi,
 \quad \epsilon\in(R/\pi R)^\times,
\]
after the common coefficient extension and compatible choices of determinant bases. The second congruence follows from the identical residual coefficient and local diagrams.

The auxiliary primes will be chosen so that
\[
 r_j\to -1,\quad r_j\text{ inert also in }L,\quad
 \rho_{M^*}(\gamma_j)\to J_j,\ J_j^2=1,\qquad
 a_{r_j}(E)\to a_{j,*}\ne0,
\]
where \(\gamma_j\) denotes sequences of rational Frobenius lifts. The CM trace \(a_{r_j}(f_0)=0\), the full CM Tate involution limit follows also by induction and determinant, and \(\gamma_j^2\) acts trivially on all ring fields of \(K\).

\begin{lemma}[CM local factorization]\label{alternation-imaginary:local-factors}
For auxiliary primes satisfying these prescriptions, put \(v_j=(1+u_j)^{\alpha_j}-1\), where \(\alpha_j\in\mathbb Z_2^\times\) is the relative inertia exponent. Replacing the full local condition by the pure singular condition on the CM problem removes a factor \(v_j^2\) from its strict determinant. Each paired measure vanishes when \(u_j=0\).
\end{lemma}
\begin{proof}
For \(M^*\) the limiting local tame/Frobenius cochains split into
\[
 U_{s,j}=[M\xrightarrow{v_j}M]\ (0,1),\qquad Q_j=[M\xrightarrow{\pm v_j} M]\ (1,2),
 \quad v_j=(1+u_j)^{\alpha_j}-1,\quad \alpha_j\in\mathbb Z_2^\times
\]
with \(M\) the scalar-extended plane locally. Use the two procyclic resolutions as in our diagram conventions (\(r_j^2\to1\)); \(\alpha_j\) uses tame generators \(\sigma_j\) generating also full relative cyclic inertia of order \(r_j+1\). At the center \(U_{s,j}\) imposes the pure singular plane as for \eqref{eq:K}, i.e. for coordinates \(s_{\rm loc},f_{\rm loc}\) by evaluation at \(\sigma_j,\gamma_j^2\) it sets \(f_{\rm loc}=0\), includes degree zero and omits degree two. Using all the \(U_{s,j}\) instead of full there on the CM problem gives
\[
                          D=\mathrm{unit}\cdot D_s\prod_j v_j^2 .        \tag{I1}\label{eq:I1}
\]
The modified complex too has square amplitude 1,2 by residual duality (pure plane orthogonality as for \eqref{eq:K}, including at residual precision as checked in the square switch). Also \(u_j\mid b^\pm\) on that problem: set \(u_j=0\), trace over the inert conductor step at high \(t=\theta\) in \eqref{eq:M2}. Each point is a descendant from omit-\(r_j\) conductor by this inert prime with transported level (canonical 2-subgroup preserved); the relative trace is Hecke with multiplier zero, and remaining weights descend. Bounded-series testing proves the assertion.
\end{proof}

\subsection{Exact central switches on the CM comparator}

We now set \(t=\mathbf u=0\) on the CM plane \(M^*\). For an initial
segment \(I\) of the auxiliary primes, impose the pure singular condition
at the primes in \(I\), and the unramified condition at the auxiliary
primes not yet switched. Write \(D_{\rm strict,I}\) for the resulting
strict determinant. Replacing the strict/full dyadic pair by the two
rational Kummer lines gives a complex \(C_{F,I}\). The determinant volumes
for this rational modification are assigned by \eqref{eq:D}, with the
logarithm basis at \(w\) and its conjugate dual.

Let \(P\) be the projected Hilbert trace and let \(Y_I\) be its descended
derivative classes from Proposition~\ref{split-local:derivative}, transported
to dual Tate by \(\lambda_A\). Use a common nonzero multiplier \(H\) for
all descent, cusp, projection, and lattice denominators, so
\(Y_\varnothing=H\lambda_A(P)\). The logarithm of a dual-Tate class here
means the logarithm after applying \(\lambda_A^{-1}\), as above. Initially,
Proposition~\ref{cm:paired} gives
\[
 D_{\rm pre}:=D_{\rm strict,\varnothing}\ne0,
 \qquad v_2(D_{\rm pre})=2v_2(e_S\log_{\omega,w}P),
 \qquad e_S=\prod_{q\in S_f}P_q(1).
\]
The following lemma identifies the two evaluations needed to preserve this
central comparison at each switch.
The finite and singular coordinates are the evaluations at
\(\gamma_j^2\) and \(\sigma_j\), respectively. The derivative identity
on the CM plane is
\[
 s_{\rm loc}(Y_{I+j})=J_jf_{\rm loc}(Y_I),\qquad
 I+j=I\cup\{j\}.
 \tag{I2}\label{eq:I2}
\]

\begin{lemma}[Exact central switch]\label{alternation-imaginary:central-switch}
Suppose the auxiliary primes satisfy the Frobenius prescriptions above.
At each step \(I\to I\cup\{j\}\), require nonzero
finite evaluations of \(Y_I\) and of a class \(z\) of opposite conjugation
sign in the problem obtained from \(C_{F,I}\) by relaxing both dyadic
conditions. Then at every step the strict complex is rationally acyclic,
\(C_{F,I}\) has cohomological ranks \((1,1)\), and \(Y_I\) spans degree
one with nonzero logarithms at both dyadic places. Moreover,
\[
 2v_2(\log_{\omega,w}Y_I)-v_2(D_{\rm strict,I})
       =2v_2(H\log_{\omega,w}P)-v_2(D_{\rm pre}).
 \tag{I3}\label{eq:I3}
\]
\end{lemma}
\begin{proof}
Proceed inductively; Proposition~\ref{alternation-imaginary:concentration}
will construct primes satisfying both evaluations. Before switching a new
prime, its unramified condition recovers the old diagram by inflation.
Initially \(H^1(C_{F,\varnothing})\) is the Heegner line with nonzero
logs by \eqref{eq:C-pair}. The finite problem has ordinary and conjugate
self-duality over the fraction field, since the fixed odd local complexes
are acyclic. These are the initial rank and duality assertions needed for
the induction.

For these derivatives one applies
\(D_I^{\rm der}=\prod_{j\in I}\sum_{a=1}^{r_j}a\sigma_j^a\) to the
upper projected point, descends its Kummer class modulo increasing
precision, and traces the Hilbert field to \(K\). Equation~\eqref{eq:K}
applies on full actual abelian Tate before projection and \(\lambda_A\):
the Frobenius-square divisibilities hold on the whole CM variety, all
CM Hecke eigenvalues here are zero, and descent uses fixed base local
ramification. This gives \eqref{eq:I2}, finite Kummer at two, pure singular
conditions on \(I\), and unramified conditions at the unswitched auxiliary
primes. Increase \(H\) on all compared classes together if further
switches require it; bounded invariants remove the descent ambiguities,
as in Proposition~\ref{split-local:derivative}.

Inductively let \(Y_I\) span \(H^1(C_{F,I})\) with nonzero dyadic logs.
The conjugation-stable conditions make this line an eigenspace. Relaxing
both dyadic conditions adds two local quotient lines, exchanged by
conjugation. Their boundary onto \(H^2(C_{F,I})\) is dual to the nonzero
finite localization. Its kernel therefore supplies a class \(z\) of the
opposite sign; the isotropy choices may be made equivariant by averaging
over the fraction field.

By hypothesis, the evaluations of \(Y_I\) and \(z\) at the next prime
are nonzero and lie in opposite \(J_j\)-eigenlines. The identity
\eqref{eq:I2} and the square switch preserve ranks one and one, with
\(Y_{I+j}\) spanning degree one. Cup reciprocity against \(z\) forces a
nonzero dyadic localization of \(Y_{I+j}\): the mixed cup at the new prime
is nonzero and every other odd contribution vanishes. Conjugation stability
of the new line then gives nonzero logarithms at both dyadic places.
Equation~\eqref{eq:D} restores strict acyclicity. In particular the final
modified determinant satisfies \(D_s(0,\mathbf0)\ne0\).

For the exact valuation, the left side of \eqref{eq:I3} is the
class-functional valuation with the assigned dyadic volumes. Their
comparison triangles are compatible with the separate odd switches;
the adjunctions needed here are over the fraction field. At the new prime
the mixed pairing is a unit times \(e(x,J_jy)\). Thus \eqref{eq:I2} and
the square switch, with the integral plane-coordinate volumes, preserve
this class-functional valuation. Its initial value is the right side
of \eqref{eq:I3}, including the common multiplier \(H\).
\end{proof}

\subsection{Residual concentration and compatible prime choices}

\begin{proposition}[Compatible primes and residual concentration]\label{alternation-imaginary:concentration}
A finite collection of auxiliary prime sequences can be chosen to satisfy all prescriptions of Lemmas~\ref{alternation-imaginary:local-factors} and~\ref{alternation-imaginary:central-switch}, to have nonzero fixed-place Frobenius exponents in one tame variable, and to satisfy
\[
 D\big|_{t=\pi=0}\ne0
\]
for both the CM comparator and the elliptic curve. The non-CM limiting traces \(a_{j,*}\) can all be required to be nonzero.
\end{proposition}
\begin{proof}
We construct the primes inductively, retaining three properties of the
chosen tuple:
\begin{enumerate}[label=\textup{(\roman*)}]
\item The first cohomology of the central CM finite problem is a line
generated by \(Y_I\), with nonzero dyadic logarithms, and its strict
determinant is nonzero.
\item The first tame variable has nonzero Frobenius exponent at one
place above each prime of \(S_f\).
\item Each chosen prime is inert in both \(K\) and \(L\), has the
required CM involution limit, and has nonzero limiting trace on \(E\).
\end{enumerate}
The initial central assertion is \eqref{eq:C-pair}. At each new prime
we impose the two central evaluations of
Lemma~\ref{alternation-imaginary:central-switch} to preserve it. After
the first prime, we also impose a residual rank-two evaluation whenever
the residual dimension is positive. We first show that these requirements
force termination, and then prove that they can all be met simultaneously.

\emph{Dimension reduction.} Choose the first prime \(r_0\) with the
nonzero fixed-place exponents in \(u_0\). At \(t=0\), consider the DVR
at \((\pi)\) of \(\mathcal O[[\mathbf u]]\). Write \(h\) for the
degree-one dimension of the full-at-added-primes strict problem on its
residue field. The central invariant and \eqref{eq:I1} make the generic
problem acyclic. Fixed odd local complexes are acyclic over this DVR:
their nonconstant unramified scalar has no invariants or dual invariants
against the fixed residual matrices, and the odd local Euler
characteristic is zero. The added inert local complexes are acyclic by
their inertia scalars; the dyadic degree-zero terms vanish by the
Frobenius test. Thus conjugate duality is perfect.

To apply \eqref{eq:H1} after inducing over \(\mathbb Q\), allow the
quadratic ramification primes with the unramified conditions in that
lemma. The coefficient is good there and the scalar unramified over
\(K\); unramified inflation recovers the same global problem. At the
acyclic odd places use the zero replacement allowed in that lemma.
Corollary~\ref{alternation-imaginary:parity} now makes \(h\) even. The
curve has the identical residual problem and hence the same dimension.

If \(h\ge2\) add a variable with evaluations of two independent strict classes of rank two in the finite plane at that new place (initial \(u_j=0\), residue generic field of the old variables with \(t=0\)). Over that field localization and duality are exact with either pure plane: the local action at this new place is trivial in the limit, and conjugate transport with pure cups and mixed pairing is as at \eqref{eq:K}. Thus changing to the pure singular plane drops \(h\) by two as in the square switch (the upper image is exactly the finite plane). This switched problem is a specialization of the problem using \(U_{s,j}\) before setting \(u_j=0\); with \(u_j\) also generic that condition and full give the same cohomology. Hence the new generic residual dimension is at most \(h-2\) by the perfect models. Even parity at each step then terminates with
\[
                              D|_{t=\pi=0}\ne0.
\]

It remains to justify the simultaneous choices used in this induction.
The first prime must supply the fixed-place exponents. Every later prime
must supply the residual rank-two evaluation when \(h>0\). At every
step the two central evaluations preserve the first invariant, and the
non-CM trace must stay nonzero. The following constructions impose these
requirements in that order without losing the earlier ones.

\emph{Ring-class characters and the first variable.} In \(\operatorname{Pic}(\mathcal O_K)\) take cyclic factor generators represented by prime ideals \(I_d\) at distinct split primes outside the support and away from use, by Chebotarev. Write \(I_d^{h_d}=(x_d)\) for their order relations. For the chosen \(v\)'s put \((y_v)=v\prod I_d^{-n_{vd}}\). At inert \(r_j\to-1\) with \(2^{m_j}\mid r_j+1\), use the \(2^{m_j}\)-power residue symbol on \(\mathbb F_{r_j^2}^\times\) with compatible roots at the chosen place; it kills rational residues and maps onto the cyclic roots of that order. Require zero exponents on \(x_d\); then by the order class sequence the residue character extends killing the \(I_d\)'s on the ring class group, with exponents at \(v\) given up to conventions by \(y_v\). This produces the specified independent conductor variables (their relative inertia groups for the different primes multiply, using units just signs).

Take \(\gamma_j\) approaching elements \(g_i=c\eta_i\) (\(i\) indexes stages) where \(\eta_i\in G_K\) fixes \(W,\mu_{2^\infty}\). Thus \(g_i\) is also in the inert \(L\)-coset. Exponents on the above numbers of \(K^\times\) are read modulo \(2^{m_j}\) from the Kummer translations by \(g_i^2\), i.e. translations by \(\eta_i\) minus their values by \(\eta_i\) on the conjugate numbers. At the first step we can prescribe zero on \(x_d\) and all differences nonzero on \(y_v\) with fixed values as required. Indeed the numbers and their conjugates are jointly multiplicatively independent by valuations, and Kummer images on the specified kernel are open in the full tuple space (use valuations after the finite extension and rational restriction injectivity on passing to the cyclotomic tower by a central nonidentity cyclotomic scalar). Take these initial \(\eta_i\) constant. At subsequent steps instead require \(\eta_i\) trivial also on all the radicals, giving zeros on the class relations again.

\emph{Residual evaluations.}
In the residual tests for those subsequent steps use the evaluation groups \(\mathcal G_K=\prod_i G_K,\ \mathcal G=\mathcal G_K\rtimes\langle c\rangle\) of the diagram conventions. Degree one in question injects into unrestricted crossed classes (in particular by dyadic invariant vanishing just checked). Take the termwise product kernel \(\mathcal H\) over \(K\) of \(W\), ring-class and cyclotomic data and the radicals just used (including conjugates). Restriction here detects the classes. Before radicals the quotient has a normal product of the \(C_3\)'s (other data abelian over \(K\)), whose diagonal \(C_3\) kills its invariants and abstract \(H^1\); and a lift fixing the cyclotomic data acts trivially by conjugation on radical translations, so these give no equivariant homomorphisms either. The conjugate induction is absolutely simple since \(W\) is absolutely simple over \(K\) and the two scalar twists differ in determinant (already by the infinite-order image using \(u_0\)). Joint evaluations on \(\mathcal H\) thus span both copies of the induction for the two classes, by abstract Shapiro and the simple evaluation test. Evaluation on \((c\eta)^2,\ \eta\in\mathcal H\) acts via \(1+c\) and projection back to the first summands, additively and still spanning both tested planes. Some \(\eta\) therefore gives rank two by the quadratic determinant test.

\emph{Nonzero traces on the elliptic curve.}
Now perturb either choice of \(g\) by products of squares within \(\mathcal H\), preserving the residual test as well as any exact radical prescriptions. We can thus ensure \(a_{j,*}\ne0\) on \(E\), since these kernels contain uniformly an open determinant-one piece on the curve's Tate data (non-CM open image with bounded-derived-length excluded data as in the evaluation discussion). Products of squares there still contain such a piece, so the coset trace is not identically zero even on limiting matrices. Preserve the nonvanishing by keeping henceforth \(E[2^a]\) actions fixed to sufficient fixed depth.

\emph{The two characteristic-zero evaluations.}
Lastly use squares within the deeper kernel fixing also \(E[2^a]\) and full CM Tate of \(A\). For the two central rational classes \(Y_I,z\) of the required opposite signs, evaluation still detects each by restriction here. Indeed these Selmer classes with finite or relaxed dyadic conditions inject into abstract crossed classes as before. A constant element \(h_0 c h_0 c\), with \(h_0\in G_{K L(E[2^a])}\) of cyclotomic value a fixed integer \(m>1\), acts by homothety \(m\) on CM Tate and \(M^*\) (diagonal CM characters with product cyclotomic, swapped by \(c\)), trivially on ring fields and the other designated finite data. It is central modulo the pre-radical kernel. And it acts on radical translations by the integer \(m^2\); so even abstract equivariant homomorphisms from the additional kernel to the plane vanish. Thus evaluations of either detected class on the deeper kernel span the simple characteristic-zero plane over \(K\). Multiplying \(g\) by squares there adjusts evaluations at \(g^2\) for a sign-\(\epsilon\) class by \(2(1+\epsilon J_j)\) applied to its kernel evaluations, additively: conjugation by \(g\) uses the eigenlaw on classes, with coboundaries trivial here. These projected changes span the required eigenline for each. Hence in characteristic zero the two nonzero tests can be met together without changing previous requirements.

Chebotarev at increasing precision now gives the primes, retaining the previous cochain evaluations and all stated field and Frobenius prescriptions, including the exact class-relation zeros to cofinal symbol depth. Local generators can be taken from inertia with the needed relative cyclic action. This validates \eqref{eq:I2}--\eqref{eq:I3} and the residual-dimension induction simultaneously (central and residual cocycles use their respective finite diagrams, extended before limits).
\end{proof}

\subsection{Leading coefficient and divisibility tests}

For the CM comparator, Lemma~\ref{alternation-imaginary:local-factors}
already gives integral canceled series
\[
 \widetilde B=\frac{B}{\prod_j u_j^2},\qquad
 \widetilde D=\frac{D}{\prod_j u_j^2}.
\]
Indeed \(v_j/u_j\) is a unit. The central switches give
\(\widetilde D(0,\mathbf0)\ne0\). We compare these two scalar values
before proving divisibility of the full series.

\begin{lemma}[The central mixed coefficient]\label{alternation-imaginary:mixed-coefficient}
For \(f_0\), with \(I\) the set of all tame indices and with the same common multiplier \(H\) as in Lemma~\ref{alternation-imaginary:central-switch}, one has
\[
             v_2\left(\left.\frac{b^\pm}{\prod u_j}\right|_{t=\mathbf u=0}\right)
                         =v_2(e_S\log_{\omega,w}Y_I/H).              \tag{I4}\label{eq:I4}
\]
Consequently both \(\widetilde B(0,\mathbf0)\) and
\(\widetilde D(0,\mathbf0)\) are nonzero, and their ratio is a dyadic
unit:
\[
 v_2\!\left(
 \frac{\widetilde B(0,\mathbf0)}{\widetilde D(0,\mathbf0)}
 \right)=0.
\]
\end{lemma}
\begin{proof}
We retain the descent multiplier throughout the coefficient calculation. At \(t=0\) use the disk-center formula of \eqref{eq:M1}--\eqref{eq:M3} with \(\Psi_{\mathbf u}\) the tame weights. Euler operators all act by translations, and opposite original-level sums differ up to sign and unit translation by Atkin--Lehner as there. The projected point at conductor \(\prod r_j\) in the sum is denoted \(X_I\).

To check the positive-orientation formula with additive congruences, first multiply the stage formula by \(H G(\Psi_{\mathbf u}(\mathrm{Fr}_w)^{-1})\), where \(G\) is a fixed Frobenius polynomial with \(G(1)\ne0\) sending all points on the actual abelian factor over dyadic unramified local fields into the small formal-log range. Use unramified Néron base change, component order, the abelian and toric Frobenius relations and a characteristic power for unipotent reduction, then a fixed integer in the formal range as in the unramified log comparison. Here \(\mathrm{Fr}_w\) uses the corresponding CM class action. Thus the weighted logs are now of translates of \(G(\mathrm{Fr}_w) H X_I\), uniformly bounded and with log congruences for division up to the same fixed loss. All coefficients with a variable dropped from the weighted sum vanish by the inert Hecke traces.

Work modulo increasing precision with each \(u_j^2=0\) (group powers map there, absorbing fixed Euler denominators). Write the labels as Hilbert representatives times \(\prod_j\sigma_j^{k_j}\). In the mixed coefficient only weights \(\prod_j\alpha_{j,i} k_j\) contribute (\(\alpha_{j,i}\) stage exponents): expanding the product of exponents of the full labels, all other terms drop at least one relative weight and use trace zero. Euler multipliers contribute just \(e_S\). The resulting additive sum without these unit exponents and Euler multipliers is the log on the Hilbert-label sum of translates of \(G(\mathrm{Fr}_w)H D_I^{\rm der}X_I\). Its limit is \(G(1)\log_\omega Y_I\). Indeed restriction back up of the traced descended class is exactly the corresponding Hilbert sum before \(G\) in Kummer modulo growing precision. This comparison takes place on full actual abelian Tate after the common multiplier and before final lattice projection/\(\lambda_A\); the descended trace classes at \(w\) are represented by local Kummers as in the descent proof. Thus the point congruence after restriction is modulo division in the upper \textbf{unramified} local fields, so transfers to a uniform additive congruence after \(G\). On the base-field local point \(G\) acts by \(G(1)\); base logs of such finite-level representatives converge to the localization log along the cohomology models at the fixed group (Kummer injection and compactness), with the stated polarization convention. Passing precisions therefore computes the mixed coefficient. On the series side multiplying by the group polynomial likewise uses only \(G(1)\) in that coefficient, by the \(u_j\)-divisibilities. Cancelling yields \eqref{eq:I4}, also at the opposite orientation (same Euler constants). More explicitly, \eqref{eq:I3} and the initial value of \(D_{\rm pre}\) give
\[
v_2(D_{\rm strict,I})
 =2v_2(\log_\omega Y_I)-2v_2(H)+2v_2(e_S).
\]
By \eqref{eq:I4}, this is the valuation of
\(\widetilde B(0,\mathbf0)\), and it is finite because the central
derivative has nonzero logarithm. Equation~\eqref{eq:I1} identifies
\(\widetilde D(0,\mathbf0)\) with \(D_{\rm strict,I}\) up to a
unit. The two scalar values therefore have the same valuation. This
proves the stated central ratio, retaining every power of two in \(H\).
\end{proof}

\begin{proposition}[Integral quotient and its unit value]\label{alternation-imaginary:unit}
Embed the arithmetic coefficient ring into the common integral measure coefficient ring. For each of the two coefficient planes, \(D\mid B\) in the resulting integral series ring, and both quotients in that ring are units.
\end{proposition}
\begin{proof}
We first establish integral divisibility, then use the central scalar
ratio to prove that the quotient is a unit. At high finite \(t=\theta\)
use \(\mathscr A=\mathcal O(\theta)[[\mathbf u]][1/2]\), where
\(\mathcal O(\theta)\) is the integer ring after extension; test at its
height-one DVRs as in \eqref{eq:B1}, \eqref{eq:D2}. We spell out the
applicability here.

\begin{enumerate}[label=\textup{(\roman*)}]
\item The split dyadic Kummer lines and weighted log identification over extended constants of the local comparisons there apply. The \(\mathbf u\)-actions are unramified with some nonzero Frobenius exponent by construction. Thus the substitution from the unramified Shapiro tower (after splitting \(\theta\) locally as the ramified-field character and unramified part) is flat as discussed there. Local Heegner fields of the combined conductor contain the fixed ramified character field and arbitrarily large fixed unramified layers eventually, since they cut \(\theta\) and have the growing unramified orders. High \(\theta\) is nontrivial on inertia after the indicated fixed reduction extensions. These are precisely the hypotheses for bounded Kummer/log transfer also on projected abelian Tate and several-variable cyclic quotients.
\item Fixed odd singular determinants removed in imposing the characteristic-zero unramified conditions multiply to \(D_{\rm odd}=\prod_{q\in S_f,\ q\ne2}P_q(\chi(\mathrm{Fr}_v))P_q(\chi(\mathrm{Fr}_v)^{-1})\) up to units, \(\chi=\theta\Psi_{\mathbf u}\), \(v\mid q\). Use the local coinvariant computation of \eqref{eq:B1} (for the CM form this is the usual good/bad local Euler factor on the induced character realization; inertia has finite image and the Hecke \(L\)-function gives the Frobenius polynomial on its unramified part). None vanish on a whole \(u_j=0\), by the all-zero tame specialization and Frobenius weights. The added inert locals are acyclic on the test except possibly on the CM branch at primes \((u_j)\). Indeed away from these divisors tame inertia alone suffices; on the curve branch the Frobenius-square limit minus one is also invertible by \(a_{j,*}\ne0\), determinant \(-1\) on the original Frobenius limit, and all the scalar weights trivial on its square.
\item Formula \eqref{eq:M2} on these conductors and the weighted-log identification give \eqref{eq:B1} in the form \(B(\theta,\mathbf u)=\mathrm{unit}\cdot p_0^2 D_{\rm odd}\) over extended 2-inverted constants. Here \(p_0\in\mathscr A\) is the localized weighted Heegner class coordinate at \(w\) on the free Kummer line, up to common point/descent multipliers (using the unnormalized class sums). Membership and log comparisons are on finite cyclic quotients before the unramified limits exactly as in the local comparison, with no division by a growing degree. In the exceptional \((u_j)\) tests on \(f_0\), \eqref{eq:I1} has order exactly two at \(u_j\) for high \(\theta\), since \(D_s(0,\mathbf0)\ne0\). The \(u_j\)-divisibilities of both \(b^\pm\) hence show \(p_0^2/(D(\theta)/D_{\rm odd})\) has no pole there (order in \(u_j\) is unchanged by extending constants).

\item At every other test, if \(B(\theta)\ne0\), apply the square switch proof of \eqref{eq:D2} with that primitive divisor \(D(\theta)/D_{\rm odd}\) (nonzero by residual regularity). All local lines/conditions and duality at old places apply by the preceding comparisons. Degree zeros vanish by Tate absolute irreducibility over \(K\) (even before scalar character, on either plane), and Selmer \(H^1\)'s inject into unrestricted crossed classes including at surplus fibers. The new derivative primes for this test are \emph{additional} inert primes, tending to conjugation on full Tate data as at \eqref{eq:K}, not changing the variables. That proof applies on the modular abelian factor also, with fixed dyadic conductor and old moving conductors unramified at the required fixed places, high-character Kummer membership as above. For each surplus test the conjugate induction on the fiber is absolutely simple by Tate simplicity and distinct twisted determinants (high dyadic inertia in \(\theta\)). Restriction to product kernels fixing full Tate and ring-class data detects evaluations by the same central homothety trivial on ring fields used at \eqref{eq:D2}; hence squares of elements \(c\eta\) with \(\eta\) there provide rank-two evaluations including one of a primitive reduction, as in that proof. Thus \eqref{eq:K} and exact plane comparisons give the nonnegative class-functional valuation at the DVR by switching to minimal fiber rank with integral classes (fixed multipliers are units). This says \(p_0^2/(D(\theta)/D_{\rm odd})\) has no pole here either.

\end{enumerate}

Consequently normality and \eqref{eq:B1} give
\(B(\theta)\in(D(\theta))\) after inverting two in the extended
series ring; this is automatic if \(B(\theta)=0\). Residual regularity
at \(t=0\), a tame Weierstrass change, and the high-character remainder
test give \(D\mid B\) integrally for both planes.

Let \(U_0=B_0/D_0\) be the integral CM quotient. Canceling the common
\(\prod_j u_j^2\) gives \(\widetilde B=U_0\widetilde D\), so its
constant coefficient is the scalar ratio computed in
Lemma~\ref{alternation-imaginary:mixed-coefficient}. It is a dyadic unit,
and hence \(U_0\) is a unit. Finally the congruences
\(B_E\equiv B_0\) and \(D_E\equiv\epsilon D_0\pmod\pi\), together
with \(D_0\bmod\pi\ne0\), identify the quotient reductions by
cancellation. Thus the elliptic quotient is also a unit.
\end{proof}

\subsection{The ordinary center and the exact product formula}

\begin{proof}[Proof of Proposition~\ref{alternation-imaginary:anchor}]

Use the unit identity of Proposition~\ref{alternation-imaginary:unit} first at \(t=0\) on a group-power tame line retaining generic nonvanishing and nonzero dyadic Frobenius exponent, localized at its characteristic-zero origin. The rank-detection argument in the proof of Proposition~\ref{split-local:center-rank} applies there, with the following local checks: dyadic conditions use the unramified Kummer and log comparisons at that DVR, specializing to usual Kummer. All odd local complexes are contractible there including the inert moving ones by \(a_{j,*}\ne0\). Thus the all-dyadic-finite problem has special rank one by \(s_2(E/K)=1\) and inflation, with only the paired lines in special cohomology. The \(t=0\) log formula gives the square of the weighted trace localization up to units there (all Euler terms nonzero at the origin by their weights); again one may first multiply by the Frobenius bound of the unramified comparison. Strict generic acyclicity and nonzero log then give the generic line assertions by \eqref{eq:D}, and zero valuation of the class-functional tensor by the unit equality. Hence the global class specializes nontrivially exactly as in the rank implication there. This specialization is the Heegner Hilbert trace on the curve times \(\prod_j a_{j,*}\) (and any fixed nonzero class multiplier), by the inert Hecke relations. By \eqref{eq:GZ-E} the curve's split product therefore has a simple zero and hence finite Tate--Shafarevich over \(K\).

At the all-zero center, let \(F_j\) be the limiting Frobenius matrix on the curve at the \(j\)-th inert prime. It has \(\det F_j=-1\) and \(\operatorname{tr}F_j=a_{j,*}\), so
\[
 \det(F_j^2-1)=\det(F_j-1)\det(F_j+1)=-a_{j,*}^2.
\]
The full-to-unramified singular block is in degrees \(1,2\), with Frobenius square and residue size tending to one. Its determinant therefore has valuation \(2v_2(a_{j,*})\). Inflation identifies the remaining strict complex with the original strict complex. On the numerator \eqref{eq:M3} the inert traces in each orientation give exactly the same valuation multipliers. Cancelling therefore gives the equality of valuations for the true strict complex and paired conductor-one log used at \eqref{eq:E}, now with support \(S_f\) (the calculation there includes extra split good primes just the same way by Haar). Thus \eqref{eq:E} and \eqref{eq:G} give \(X(E)+X(E^k)=0\), with the general split odd discriminant \(k<-4\); here the split Gross--Zagier normalization and isogeny comparison are as specified at \eqref{eq:GZ-E}. Restoring the original notation, this proves
\[
 \an(E^h)+\an(E^{hk})=1,\qquad X(E^h)+X(E^{hk})=0.
\]
The product \(h\) is an odd positive fundamental discriminant, or \(1\), prime to \(2N_E\); the odd negative discriminant \(k\) is prime to \(2hN_E\), and every prime dividing \(2hN_E\) splits in \(K\). These are precisely the requirements of Proposition~\ref{prop:anchor}.
\end{proof} 
